分别表示有交互效应和无交互效应时,期望响应作为因子水平的函数的图形.无交互效应意味着对一切 $l$ 与 $k$,有 $\gamma_{lk}=0$.

\begin{mvfigure}
\begin{tikzpicture}[x=1.15cm,y=1.45cm,>=Stealth,font=\zihao{6}]
  \draw[->] (0,0) -- (4.75,0);
  \draw[->] (0,0) -- (0,2.15);
  \node[below] at (2.3,-0.52) {因子2的水平};
  \node[rotate=90] at (-0.30,1.05) {期望响应};
  \foreach \x in {1,2,3,4}{%
    \draw (\x,0) -- (\x,-0.07);
    \node[below] at (\x,-0.09) {\x};}
  \draw (1,0.70)--(2,1.50)--(3,0.70)--(4,1.65);
  \draw (1,1.00)--(2,1.20)--(3,1.00)--(4,1.10);
  \draw (1,0.40)--(2,1.00)--(3,0.35)--(4,0.80);
  \node[anchor=west] at (4.05,1.80) {因子1的水平1};
  \node[anchor=west] at (4.05,1.45) {因子1的水平3};
  \node[anchor=west] at (4.05,1.10) {因子1的水平2};
  \node at (2.3,-1.05) {(a)};
\end{tikzpicture}

\begin{tikzpicture}[x=1.15cm,y=0.62cm,>=Stealth,font=\zihao{6}]
  \draw[->] (0,0) -- (4.75,0);
  \draw[->] (0,0) -- (0,4.55);
  \node[below] at (2.3,-1.25) {因子2的水平};
  \node[rotate=90] at (-0.30,2.25) {期望响应};
  \foreach \x in {1,2,3,4}{%
    \draw (\x,0) -- (\x,-0.16);
    \node[below] at (\x,-0.20) {\x};}
  \draw (1,2.60)--(2,3.40)--(3,2.60)--(4,4.00);
  \draw (1,2.00)--(2,2.80)--(3,2.00)--(4,3.40);
  \draw (1,1.40)--(2,2.20)--(3,1.40)--(4,2.80);
  \node[anchor=west] at (4.05,4.25) {因子1的水平3};
  \node[anchor=west] at (4.05,3.65) {因子1的水平1};
  \node[anchor=west] at (4.05,3.05) {因子1的水平2};
  \node at (2.3,-2.35) {(b)};
\end{tikzpicture}
\phantomsection\label{fig:6.3}\mvcap{图 6.3\quad 期望响应曲线：(a) 有交互效应；(b) 无交互效应}
\end{mvfigure}

类似于\hyperref[eq:6-55]{式(6-55)},每个观测值也可分解为
\phantomsection\label{eq:6-56}
\begin{equation}
x_{lkr} = \bar{x} + (\bar{x}_{l\cdot}-\bar{x}) + (\bar{x}_{\cdot k}-\bar{x}) + (\bar{x}_{lk}-\bar{x}_{l\cdot}-\bar{x}_{\cdot k}+\bar{x}) + (x_{lkr}-\bar{x}_{lk}) \tag{6-56}
\end{equation}

其中 $\bar{x}$ 为总平均, $\bar{x}_{l\cdot}$ 为因子 1 的第 $l$ 个水平下观测值的平均, $\bar{x}_{\cdot k}$ 为因子 2 的第 $k$ 个水平下观测值的平均, $\bar{x}_{lk}$ 为因子 1 的第 $l$ 个水平、因子 2 的第 $k$ 个水平下观测值的平均.将 $(x_{lkr}-\bar{x})$ 取平方再求和,可得

\phantomsection\label{eq:6-57}
\begin{equation}
\begin{aligned}
\sum_{l=1}^{g}\sum_{k=1}^{b}\sum_{r=1}^{n}(x_{lkr}-\bar{x})^2
&= \sum_{l=1}^{g} bn(\bar{x}_{l\cdot}-\bar{x})^2 + \sum_{k=1}^{b} gn(\bar{x}_{\cdot k}-\bar{x})^2 \\
&\quad + \sum_{l=1}^{g}\sum_{k=1}^{b} n(\bar{x}_{lk}-\bar{x}_{l\cdot}-\bar{x}_{\cdot k}+\bar{x})^2 \\
&\quad + \sum_{l=1}^{g}\sum_{k=1}^{b}\sum_{r=1}^{n}(x_{lkr}-\bar{x}_{lk})^2
\end{aligned}
\tag{6-57}
\end{equation}

或 \quad $SS_{总}=SS_{因子1}+SS_{因子2}+SS_{交互}+SS_{残差}$

与\hyperref[eq:6-57]{式(6-57)}中各平方和相联系的相应自由度为:
\phantomsection\label{eq:6-58}
\begin{equation}
gbn-1 = (g-1)+(b-1)+(g-1)(b-1)+gb(n-1) \tag{6-58}
\end{equation}

ANOVA 表有如下形式:

\begin{mvtable}
\mvtabcap{比较两因子效应及其交互效应的方差分析表}
\begin{tabular}{@{}ccc@{}}
\toprule
变差来源 & 平方和 (SS) & 自由度 (d.f.)\\
\midrule
因子 1 & $SS_{因子1}=\sum\limits_{l=1}^{g} bn(\bar{x}_{l\cdot}-\bar{x})^2$ & $g-1$\\[8pt]
因子 2 & $SS_{因子2}=\sum\limits_{k=1}^{b} gn(\bar{x}_{\cdot k}-\bar{x})^2$ & $b-1$\\
\bottomrule
\end{tabular}
\end{mvtable}
